\begin{figure}[ht]
		\centering
		\begin{tikzpicture}[>=Stealth, line width=0.8pt]
			
			\def\L{4.0}
			\def\a{1.5}
			\def\F{0.6}
			\def\yrea{0.55} % altezza frecce reazioni
			
			% ============================================================
			% PANNELLO (a) — sistema originale con carico
			% ============================================================
			\begin{scope}[xshift=0cm]
				
				% trave
				
				\draw[line width=1.5pt] (0,0) -- (\L,0);
				
				% cerniera in A
				\draw[thick] (0,0) -- (-0.22,-0.38) -- (0.22,-0.38) -- cycle;
				\draw[thick] (-0.32,-0.38) -- (0.32,-0.38);
				\foreach \x in {-0.24,-0.12,0,0.12,0.24}{
					\draw[thin] (\x,-0.38) -- (\x-0.10,-0.52);
				}
				
				
				% carrello verticale in C
				\draw[thick] (\L/2,0) -- (\L/2-0.22,-0.38)
				-- (\L/2+0.22,-0.38) -- cycle;
				\draw[thick] (\L/2-0.15,-0.44) circle (0.07);
				\draw[thick] (\L/2+0.15,-0.44) circle (0.07);
				\draw[thick] (\L/2-0.32,-0.52) -- (\L/2+0.32,-0.52);
				\foreach \x in {-0.24,-0.12,0,0.12,0.24}{
					\draw[thin] (\L/2+\x,-0.52) -- (\L/2+\x-0.10,-0.66);
				}
				
				% carrello verticale in B
				\draw[thick] (\L,0) -- (\L-0.22,-0.38)
				-- (\L+0.22,-0.38) -- cycle;
				\draw[thick] (\L-0.15,-0.44) circle (0.07);
				\draw[thick] (\L+0.15,-0.44) circle (0.07);
				\draw[thick] (\L-0.32,-0.52) -- (\L+0.32,-0.52);
				\foreach \x in {-0.24,-0.12,0,0.12,0.24}{
					\draw[thin] (\L+\x,-0.52) -- (\L+\x-0.10,-0.66);
				}
				
				% punti notevoli
				\node[circle, fill=white, draw=black, inner sep=0pt,
				minimum size=5pt, line width=0.8pt] at (0,0) {};
				\node[above, font=\scriptsize] at (0,0) {$A$};
				
				\node[circle, fill=white, draw=black, inner sep=0pt,
				minimum size=5pt, line width=0.8pt] at (\L,0) {};
				\node[above, font=\scriptsize] at (\L,0) {$B$};
				
				% mezzaluna in C (vincolo attivo)
				\fill[white] (\L/2-0.08,0) arc(180:360:0.08) -- cycle;
				\draw[thick]  (\L/2-0.08,0) arc(180:360:0.08);
				
				
				\node[above, font=\scriptsize] at (\L/2,0) {$C$};
				\node[below, font=\scriptsize] at (\a,0)   {$P$};
				
				
				% forza F in P
				\draw[->, thick] (\a, \F) -- (\a, 0)
				node[midway, left, font=\scriptsize] {$F$};
				
				% quota a
				\draw[<->, thin] (0,-0.9) -- (\a,-0.9)
				node[midway, above, font=\scriptsize] {$a$};
				
				% quota ell
				\draw[<->, thin] (0,-1.3) -- (\L,-1.3)
				node[midway, above, font=\scriptsize] {$\ell$};
				
				\node[font=\scriptsize] at (\L/2, -1.8) {(a)};
				
			\end{scope}
			
			% ============================================================
			% PANNELLO (b) — sistema principale e stato di autoequilibrio
			% ============================================================
			\begin{scope}[xshift=6.5cm]
				
				% trave
				\draw[line width=1.5pt] (0,0) -- (\L,0);
				
				% cerniera in A
				\draw[thick] (0,0) -- (-0.22,-0.38) -- (0.22,-0.38) -- cycle;
				\draw[thick] (-0.32,-0.38) -- (0.32,-0.38);
				\foreach \x in {-0.24,-0.12,0,0.12,0.24}{
					\draw[thin] (\x,-0.38) -- (\x-0.10,-0.52);
				}
				
				% carrello verticale in C (grigio — ridondante)
				\draw[gray!50, thick] (\L/2,0) -- (\L/2-0.22,-0.38)
				-- (\L/2+0.22,-0.38) -- cycle;
				\draw[gray!50, thick] (\L/2-0.15,-0.44) circle (0.07);
				\draw[gray!50, thick] (\L/2+0.15,-0.44) circle (0.07);
				\draw[gray!50, thick] (\L/2-0.32,-0.52)
				-- (\L/2+0.32,-0.52);
				\foreach \x in {-0.24,-0.12,0,0.12,0.24}{
					\draw[gray!50, thin] (\L/2+\x,-0.52)
					-- (\L/2+\x-0.10,-0.66);
				}
				
				
				% carrello verticale in B
				\draw[thick] (\L,0) -- (\L-0.22,-0.38)
				-- (\L+0.22,-0.38) -- cycle;
				\draw[thick] (\L-0.15,-0.44) circle (0.07);
				\draw[thick] (\L+0.15,-0.44) circle (0.07);
				\draw[thick] (\L-0.32,-0.52) -- (\L+0.32,-0.52);
				\foreach \x in {-0.24,-0.12,0,0.12,0.24}{
					\draw[thin] (\L+\x,-0.52) -- (\L+\x-0.10,-0.66);
				}
				
				
				% punti notevoli
				\node[circle, fill=white, draw=black, inner sep=0pt,
				minimum size=5pt, line width=0.8pt] at (0,0) {};
				\node[above, font=\scriptsize] at (0,0) {$A$};
				
				\node[circle, fill=white, draw=black, inner sep=0pt,
				minimum size=5pt, line width=0.8pt] at (\L,0) {};
				\node[above, font=\scriptsize] at (\L,0) {$B$};
				
				%				% mezzaluna in C (vincolo ridondante, semiluna grigia)
				\fill[white]          (\L/2-0.08,0) arc(180:360:0.08) -- cycle;
				\draw[gray!50, thick] (\L/2-0.08,0) arc(180:360:0.08);
				
				
				\node[above, font=\scriptsize] at (\L/2,0) {$C$};
				
				
				% reazioni stato di autoequilibrio
				% Y_A = -1/2 (verso il basso): freccia traslata sopra la trave
				\draw[->, thick] (0, 2*\yrea) -- (0, \yrea)
				node[midway, left, font=\scriptsize] {$\tfrac{1}{2}$};
				% Y_C = 1 (verso l'alto): freccia traslata sopra la trave
				\draw[->, thick] (\L/2, \yrea) -- (\L/2, 2*\yrea)
				node[midway, right, font=\scriptsize] {$1$};
				% Y_B = -1/2 (verso il basso): freccia traslata sopra la trave
				\draw[->, thick] (\L, 2*\yrea) -- (\L, \yrea)
				node[midway, right, font=\scriptsize] {$\tfrac{1}{2}$};
				
				\node[font=\scriptsize] at (\L/2, -1.8) {(b)};
			\end{scope}
			
		\end{tikzpicture}
		\caption{Trave iperstatica di grado 1
			(\ESEref{ex:iper_stat_1}):
			\emph{(a)}~sistema originale con carico $F$ in $P$;
			\emph{(b)}~sistema principale (carrello in $C$ in grigio)
			e stato di autoequilibrio $\bm{r}_1$: reazione unitaria
			verso l'alto in $C$, reazioni $1/2$ verso il basso in $A$
			e $B$.}
		\label{fig:iper_stat_1_dati}
	\end{figure}
